\section{幺正变换与场算符}

Gross 第 4 章：二次量子化中的基变换是 Fock 空间上的幺正变换；场算符把坐标表象与占据数表象连起来。

\subsection{单粒子基变换}

若新轨道 $\phi'_\beta(\mathbf{r})=\sum_\alpha\phi_\alpha(\mathbf{r})U_{\alpha\beta}$，$U^\dagger U=1$，则产生湮灭算符
\begin{equation}
  c'_\beta=\sum_\alpha U_{\alpha\beta}^* c_\alpha,
  \qquad
  \{c'_\beta,(c'_\gamma)^\dagger\}=\delta_{\beta\gamma}.
\end{equation}
存在 Fock 空间幺正算符 $\mathcal{U}$（由 $U$ 诱导），使 $c'=\mathcal{U}^\dagger c\,\mathcal{U}$。哈密顿量的矩阵元在新基下用 $c'$ 重写，物理不变。

\subsection{Bogoliubov 变换}

更一般的线性正则变换混合湮灭与产生，
\begin{equation}
  \alpha_{\mathbf{k}}
  =u_{\mathbf{k}}c_{\mathbf{k}}+v_{\mathbf{k}}c_{-\mathbf{k}}^\dagger,
  \qquad
  |u|^2-|v|^2=1
  \text{（玻色）或 }|u|^2+|v|^2=1\text{（费米）},
\end{equation}
仍保持对易/反对易关系，对应某幺正 $\mathcal{U}$，但一般\textbf{不}守恒粒子数。BEC 的位移变换 $b_0\to\sqrt{N_0}+b_0$ 与 BCS 的 $u,v$ 即此例。

\subsection{场算符}

\begin{equation}
  \psi_\sigma(\mathbf{r})=\sum_\alpha\phi_\alpha(\mathbf{r})c_{\alpha\sigma},
  \qquad
  \{\psi_\sigma(\mathbf{r}),\psi_{\sigma'}^\dagger(\mathbf{r}')\}=\delta_{\sigma\sigma'}\delta(\mathbf{r}-\mathbf{r}').
\end{equation}
任意单体算符
\begin{equation}
  A^{(1)}
  =\sum_{ij}
  \langle i|a|j\rangle
  c_i^\dagger c_j
  =\sum_\sigma\int\mathrm{d}\mathbf{r}\,
  \psi_\sigma^\dagger(\mathbf{r})
  a(\mathbf{r})
  \psi_\sigma(\mathbf{r}),
\end{equation}
与基选择无关。密度 $\hat n(\mathbf{r})=\sum_\sigma\psi_\sigma^\dagger\psi_\sigma$ 的 Fourier 分量即 Gross 第 6 章的密度算符，是等离子体、线性响应与 $S(\mathbf{q})$ 的共同起点。

\begin{note}
“二次量子化与表象无关”：先选便于求解 $H_0$ 的基（平面波、布洛赫、分子轨道），再用场算符写相互作用；图技术中的动量标签即来自该选择。
\end{note}
